\documentclass[11pt]{article}
\usepackage[margin=1in]{geometry}
\usepackage{amsmath,amssymb}
\usepackage[T1]{fontenc}
\usepackage[utf8]{inputenc}
\usepackage{lmodern}
\usepackage{hyperref}
\usepackage{enumitem}

\title{Align Pitch Observability: Current Model and Proposed Direction}
\author{Codex}
\date{\today}

\begin{document}
\maketitle

\section{Problem}
On real log \texttt{ubx\_raw\_20260328\_153757.bin}, the align filter shows a persistent down/pitch-like error of about \(3^\circ\). The current filter can initialize and refine yaw reasonably, but pitch has very weak online observability after bootstrap. The result is:
\begin{itemize}[leftmargin=1.5em]
\item pitch/down error can persist for long periods,
\item gravity updates are too sparse to remove it,
\item direct longitudinal scalar pitch corrections are unstable.
\end{itemize}

\section{Notation}
\begin{itemize}[leftmargin=1.5em]
\item \(b\): body / IMU sensor frame.
\item \(v\): vehicle frame.
\item \(q_{vb}\): quaternion mapping vehicle frame to body frame.
\item \(C_{bv}(q_{vb})\): rotation matrix mapping vehicle-frame vectors into body frame.
\item \(f_b\): accelerometer specific-force measurement in body frame, i.e. the body-frame accelerometer vector used by align.
\item \(a_{\text{long}}\): longitudinal acceleration inferred from GNSS velocity change in vehicle coordinates.
\item \(a_{\text{lat}}\): lateral acceleration inferred from GNSS velocity change in vehicle coordinates.
\item \(g\): gravity magnitude.
\item \(\delta\phi,\delta\theta_p,\delta\psi\): small roll, pitch, and yaw misalignment errors.
\item \(b_a\): accelerometer bias term.
\end{itemize}

\section{Current Align State}
The align state is a small mount-misalignment error state
\[
\delta \theta =
\begin{bmatrix}
\delta \phi & \delta \theta_p & \delta \psi
\end{bmatrix}^T
\]
with covariance \(P \in \mathbb{R}^{3 \times 3}\), where:
\begin{itemize}[leftmargin=1.5em]
\item \(\delta \phi\): roll misalignment,
\item \(\delta \theta_p\): pitch misalignment,
\item \(\delta \psi\): yaw misalignment.
\end{itemize}

The nominal mount quaternion is \(q_{vb}\), mapping vehicle frame \(v\) to body/IMU frame \(b\).

\section{Current Observation Model}
\subsection{Stationary gravity}
During stationary windows, the filter low-pass filters gravity in body frame and applies gravity observations on the accelerometer channels:
\[
f_b \approx C_{bv}(q_{vb}) \begin{bmatrix} 0 \\ 0 \\ -g \end{bmatrix}
\]
with yaw masked out. This gives roll/pitch observability from gravity.

This is mathematically sound, but only when the vehicle is quasi-static and the gravity direction is representative of the vehicle frame.

\subsection{Dynamic horizontal motion}
For each GNSS interval, the code computes GNSS horizontal acceleration in vehicle coordinates:
\[
a_n = \frac{v_{n,k} - v_{n,k-1}}{\Delta t},
\quad
a_{\text{long}} = \hat t^T a_n,
\quad
a_{\text{lat}} = \hat \ell^T a_n
\]
where \(\hat t\) is the unit tangent from midpoint GNSS velocity and \(\hat \ell\) is the lateral unit vector.

It also computes a body-side horizontal specific-force vector by removing gravity projection:
\[
f^{\text{horiz}}_b
=
\bar f_b - (\bar f_b^T \hat g_b)\hat g_b
\]
then rotates that into vehicle frame:
\[
f^{\text{horiz}}_v = C_{vb}(q_{vb})^T f^{\text{horiz}}_b
\]

The current dynamic residual is then reduced to a signed horizontal heading error:
\[
e_\psi
=
\operatorname{atan2}
\left(
f^{\text{horiz}}_{v,x} a_{\text{lat}} - f^{\text{horiz}}_{v,y} a_{\text{long}},
\;
f^{\text{horiz}}_{v,x} a_{\text{long}} + f^{\text{horiz}}_{v,y} a_{\text{lat}}
\right)
\]
and applied as a yaw-only correction:
\[
\delta \psi \leftarrow \delta \psi - K_\psi e_\psi
\]

\subsection{Turn gyro}
During turn-valid windows, gyro residuals provide some additional tilt support. In practice this contributes a small amount to pitch, but not enough to remove the persistent \(3^\circ\) offset.

\section{Why Pitch is Weakly Observable in the Current Design}
The current dynamic branch is designed primarily for yaw:
\begin{itemize}[leftmargin=1.5em]
\item gravity updates affect roll/pitch only,
\item horizontal dynamic updates affect yaw only,
\item turn-gyro gives a small secondary tilt correction.
\end{itemize}

So after bootstrap, pitch has no real drive-time correction path except:
\begin{itemize}[leftmargin=1.5em]
\item sparse stationary gravity updates,
\item small turn-gyro effects.
\end{itemize}

That matches the observed diagnostics on \texttt{153757}: horizontal updates contribute essentially zero pitch correction, while gravity contributes very little total pitch correction over the whole drive.

\section{Why the Simple Scalar Pitch Fix Failed}
The attempted scalar idea was effectively:
\[
r_{\text{long}} = a_{\text{long,gnss}} - a_{\text{long,imu}}
\]
with a small-angle interpretation
\[
r_{\text{long}} \approx g \, \delta \theta_p + b_{a,x}
\]

If that model were valid and the sign were consistent, then a slow per-window correction and a batch estimate would converge to the same answer.

But on the real log, the residual is not consistent:
\begin{itemize}[leftmargin=1.5em]
\item straight accel windows have one mean sign,
\item straight brake windows have the opposite mean sign,
\item a single constant longitudinal bias does not fit the data.
\end{itemize}

So the residual is not a stable signed measurement of pitch. It is maneuver-dependent. Because of that:
\begin{itemize}[leftmargin=1.5em]
\item per-window updates chase inconsistent residuals,
\item batching only averages inconsistent residuals,
\item lowering gain only slows the wrong result.
\end{itemize}

\section{Proper Mathematical Direction}
The correct way to introduce dynamic pitch observability is to use the full specific-force vector model, not a scalar heuristic.

\subsection{Vector specific-force model}
For straight longitudinal motion, the ideal vehicle-frame specific force is approximately
\[
f_v^{\star}
\approx
\begin{bmatrix}
a_{\text{long}} \\
0 \\
-g
\end{bmatrix}
\]
and the measured body-frame force is
\[
f_b
\approx
C_{bv}(q_{vb}) f_v^{\star} + b_a + \eta
\]
where \(b_a\) is accelerometer bias and \(\eta\) is noise/model error.

This is the right measurement equation because pitch changes the \emph{vector geometry}, not just the scalar longitudinal component.

\subsection{Linearized residual}
Define
\[
r_f = f_b - C_{bv}(q_{vb}) f_v^{\star}
\]
Then linearize with respect to small misalignment:
\[
r_f \approx H_f \, \delta\theta + b_a + \eta
\]
where \(H_f\) comes from the skew-symmetric Jacobian of the rotated vector.

For straight-line windows, the most useful component is the part of \(H_f\) that makes pitch observable. This should be used with:
\begin{itemize}[leftmargin=1.5em]
\item strong gating on straightness and low lateral contamination,
\item explicit handling of maneuver phase,
\item either an estimated bias term or a nuisance-robust weighting scheme.
\end{itemize}

\section{Recommended Estimator Structure}
The next version should not do ``one window \(\rightarrow\) one pitch correction.'' Instead:

\subsection{1. Build a buffered straight-motion pitch estimator}
Accumulate a buffer of straight accel/brake windows and solve a small weighted least-squares problem over the vector residual:
\[
\min_{\delta \theta_p}
\sum_{i \in \mathcal{S}}
w_i
\left\|
f_{b,i} - C_{bv}(q_{vb} \oplus \delta \theta_p)
\begin{bmatrix}
a_{\text{long},i} \\
0 \\
-g
\end{bmatrix}
\right\|^2
\]
or its linearized equivalent.

\subsection{2. Use conservative gating}
Only include windows with:
\begin{itemize}[leftmargin=1.5em]
\item low course rate,
\item low lateral acceleration,
\item sufficient speed,
\item sufficient longitudinal acceleration magnitude,
\item stable sign/consistency over the buffer.
\end{itemize}

\subsection{3. Feed only a small bounded correction into align}
Once the buffered estimate is consistent, inject only a small pitch correction step into the align state. This preserves smoothness and avoids sharp maneuver-to-maneuver flips.

\subsection{4. Keep yaw and pitch logically separated}
Yaw should still be primarily turn-observed. Pitch should be straight-motion observed. The current formulation mixes dynamic cues in a way that is good for yaw but too weak for pitch.

\section{Summary}
The current align formulation is mathematically suited for:
\begin{itemize}[leftmargin=1.5em]
\item gravity-based roll/pitch initialization and maintenance,
\item dynamic yaw correction from horizontal motion,
\item small additional tilt support from turn gyro.
\end{itemize}

It is \emph{not} designed to provide strong dynamic pitch correction during normal driving.

The proper fix is not to increase gain on the current scalar longitudinal residual. The proper fix is to add a new vector-based straight-motion pitch observation model, buffered and heavily gated, so that pitch becomes observable without destabilizing the filter.

\end{document}
